\chapter{Классификация существующих скалярных носителей двух источников}
\label{ch:v8-extra-mass-two-source-scalar-carriers}

\section{Активные скалярные инварианты выровненного сектора}

В текущем полнографовом parent уже присутствуют поля
\begin{equation}
 B,M\in M_{2\times3}(\mathbb C),\qquad
 z=(z_Y,z_u,z_d)\in\mathbb C^3,
 \label{eq:v8-two-source-carriers-fields}
\end{equation}
и нулевой представитель
\begin{equation}
 B_0=\begin{pmatrix}\sqrt3&0&0\\0&0&0\end{pmatrix},
 \qquad
 M_0=\begin{pmatrix}0&1&0\\0&0&1\end{pmatrix},
 \qquad z_0=0.
 \label{eq:v8-two-source-carriers-vacuum}
\end{equation}
Два простейших gauge- и channel-инвариантных радиуса равны
\begin{equation}
 T_B=\Tr(BB^\dagger),\qquad T_M=\Tr(MM^\dagger),
 \qquad (T_B,T_M)|_0=(3,2).
 \label{eq:v8-two-source-carriers-radii}
\end{equation}
Инвариантные следы и определители являются стандартными генераторами
алгебр инвариантов матричных представлений колчанов
\cite{v8-two-source-carriers-lopatin}. Их наличие, однако, ещё не задаёт
смешанный член с новым центральным гамильтонианом.

\section{Независимость двух радиальных мод}

На точном радиальном срезе
\begin{equation}
 B(r)=rB_0,\qquad M(s)=sM_0
 \label{eq:v8-two-source-carriers-radial-slice}
\end{equation}
имеем
\begin{equation}
 T_B(r)=3r^2,\qquad T_M(s)=2s^2,
 \qquad
 J_{BM}=\left.\frac{\partial(T_B,T_M)}{\partial(r,s)}\right|_{1,1}
 =\begin{pmatrix}6&0\\0&4\end{pmatrix}.
 \label{eq:v8-two-source-carriers-jacobian}
\end{equation}
Следовательно,
\begin{equation}
 \rank J_{BM}=2,\qquad \det J_{BM}=24.
 \label{eq:v8-two-source-carriers-rank}
\end{equation}
Это два независимых активных скалярных режима, а не два имени одной и той
же радиальной флуктуации.

Остальные естественные инварианты новых полей не увеличивают линейный ранг.
На том же срезе
\begin{equation}
 \det(BB^\dagger)=0,
 \qquad \det(MM^\dagger)=s^4,
 \qquad \|MB^\dagger\|_F^2=0.
 \label{eq:v8-two-source-carriers-dependent-invariants}
\end{equation}
Градиент $\det(MM^\dagger)$ при $s=1$ пропорционален градиенту $T_M$.
Квадраты spectator-полей также не дают линейного источника:
\begin{equation}
 \left.\nabla_z(|z_Y|^2,|z_u|^2,|z_d|^2)\right|_{z=0}=0.
 \label{eq:v8-two-source-carriers-spectator-zero}
\end{equation}

\section{Портальная матрица и отсутствие канонической пары}

Общий линейный портал двух радиусов к $h^0=aA+bB$ имеет вид
\begin{equation}
 V_{\rm portal}=-
 \begin{pmatrix}T_B&T_M\end{pmatrix}
 \begin{pmatrix}c_{BA}&c_{BB}\\c_{MA}&c_{MB}\end{pmatrix}
 \begin{pmatrix}a\\b\end{pmatrix}.
 \label{eq:v8-two-source-carriers-portal-matrix}
\end{equation}
В вакууме он создаёт
\begin{equation}
 \begin{pmatrix}j_A\\j_B\end{pmatrix}
 =\begin{pmatrix}3&0&2&0\\0&3&0&2\end{pmatrix}
 \begin{pmatrix}c_{BA}\\c_{BB}\\c_{MA}\\c_{MB}\end{pmatrix}.
 \label{eq:v8-two-source-carriers-coefficient-map}
\end{equation}
Эта карта имеет ранг два и nullity два. Поэтому статическая пара источников
может быть реализована даже одним ненулевым радиусом при двух свободных
портальных коэффициентах. Два carrier-поля существенны не для счёта
статических чисел, а для двумерного динамического отклика.

Диагональный выбор
\begin{equation}
 C_{\rm diag}=\operatorname{diag}(\kappa_B,\kappa_M),
 \qquad (j_A,j_B)=(3\kappa_B,2\kappa_M),
 \qquad \det\frac{\partial(j_A,j_B)}
 {\partial(\kappa_B,\kappa_M)}=6
 \label{eq:v8-two-source-carriers-diagonal-convention}
\end{equation}
удобен, но не каноничен. Обе области являются тривиальными вещественными
представлениями всех сохранённых симметрий, поэтому
\begin{equation}
 \operatorname{Hom}_{G}(\mathbb R^2_{\{T_B,T_M\}},
 \mathbb R^2_{\{A,B\}})=M_2(\mathbb R),
 \qquad \dim_{\mathbb R}=4.
 \label{eq:v8-two-source-carriers-hom-space}
\end{equation}
Симметрия не выделяет диагональ, знак, нормировку или ранг портала.
Спектральное действие может связывать скалярные блоки только через явно
заданный общий оператор и его конечную геометрию
\cite{v8-two-source-carriers-stephan,v8-two-source-carriers-krajewski}.

Непосредственная проверка унаследованного функционала даёт
\begin{equation}
 \left.\frac{\partial^2S_{\rm inherited}}
 {\partial(T_B,T_M)\,\partial(a,b)}\right|_0=0_{2\times2}.
 \label{eq:v8-two-source-carriers-inherited-zero}
\end{equation}
Итоговые реестры:
\begin{equation}
 N_{\rm carrier\ classification}=6/6,
 \qquad N_{\rm canonical\ pairing}=0/5,
 \qquad N_{\rm inherited\ portal}=0/4.
 \label{eq:v8-two-source-carriers-ledgers}
\end{equation}

\begin{theorem}[Два независимых существующих скалярных режима]
Радиусы $T_B$ и $T_M$ являются активными gauge-singlet инвариантами с
ненулевыми вакуумными значениями $3$ и $2$. Их радиальный якобиан имеет
определитель $24$, поэтому они дают два независимых динамических скалярных
режима. Все прочие проверенные естественные инварианты либо нулевые в
первом порядке, либо линейно зависимы от этих двух.
\end{theorem}

\begin{nogocriterion}[Носитель не выбирает портальную матрицу]
Нельзя отождествлять два ненулевых радиуса с двумя выведенными источниками.
Пространство допустимых порталов четырёхмерно, унаследованный смешанный
гессиан нулевой, а диагональное сопоставление $T_B\mapsto A$,
$T_M\mapsto B$ является дополнительной конвенцией.
\end{nogocriterion}

\begin{protocol}\begin{enumerate}
\item активные ненулевые радиусы: \textbf{$T_B=3$, $T_M=2$};
\item determinant радиального якобиана: \textbf{$24$};
\item независимых динамических скалярных режимов: \textbf{два};
\item determinant $MM^\dagger$ даёт новый линейный режим: \textbf{нет};
\item spectator-радиусы дают линейный отклик при $z=0$: \textbf{нет};
\item размерность допустимых портальных матриц: \textbf{$4$};
\item диагональный портал каноничен: \textbf{нет};
\item унаследованный portal-Hessian: \textbf{$0_{2\times2}$};
\item carrier-classification: \textbf{$6/6$};
\item canonical-pairing: \textbf{$0/5$};
\item inherited-portal: \textbf{$0/4$};
\item следующий гейт: \textbf{parent-origin полной портальной матрицы}.
\end{enumerate}\end{protocol}

Машинный сертификат сохранён в
\path{s2t_v8_baryon_c0_singlet_triplet_central_gap_isotypic_channel_extra_edge_mass_two_source_existing_scalar_carrier_classification_gate_results.json}.

\begin{thebibliography}{9}
\bibitem{v8-two-source-carriers-lopatin} A.~A.~Lopatin,
\emph{Invariants of Quivers under the Action of Classical Groups},
arXiv:math/0608750.
\bibitem{v8-two-source-carriers-stephan} C.~A.~Stephan,
\emph{New Scalar Fields in Noncommutative Geometry}, arXiv:0901.4676.
\bibitem{v8-two-source-carriers-krajewski} T.~Krajewski,
\emph{Constraints on the Scalar Potential from Spectral Action Principle},
arXiv:hep-th/9803199.
\end{thebibliography}